\documentclass[10pt,a4paper]{article}
\usepackage[margin=2.2cm]{geometry}
\usepackage{fontspec}
\setmonofont{DejaVu Sans Mono}[Scale=0.84]
\newfontfamily\unifont{DejaVu Serif}[Scale=0.9]
\usepackage{amsmath,amssymb,mathtools}
\setcounter{MaxMatrixCols}{20}
\usepackage{longtable,booktabs,array}
\usepackage{fancyvrb}
\usepackage[dvipsnames]{xcolor}
\usepackage[colorlinks=true,linkcolor=NavyBlue,urlcolor=NavyBlue]{hyperref}
\usepackage{microtype}
\usepackage{needspace}
\usepackage{tocloft}\setlength{\cftbeforesecskip}{0.45em}
\setlength{\parindent}{0pt}\setlength{\parskip}{0.55em}
\setlength{\emergencystretch}{3em}
\renewcommand{\arraystretch}{1.15}
\allowdisplaybreaks
\begin{document}
\title{7+1: provenance\\[0.4em]\large Sources, derivations, corrections and the verification record of the notebook 7+1.nb and of its documentation}
\author{Repository: github.com/once-ere/7+1 — prepared for Nash\\ Notebook: 7+1.nb (Wolfram Language 15.0.1)\\ Generated by tools/make\_docs.py on 2026-10-26}
\date{}
\maketitle
\begin{abstract}
This record states, for every object of 7+1.nb, where it comes from (the task statement, the notebook 4+4.nb, or the literature), the equation that defines it, how it was derived, and how it was checked. It ends with the full list of the 128 checks and their results, the independent Python cross-checks, and the hashes of the files.
\end{abstract}
\tableofcontents
\Needspace{12\baselineskip}\section*{1. Sources}\addcontentsline{toc}{section}{1. Sources}
\Needspace{8\baselineskip}\subsection*{1.1 The task statement}\addcontentsline{toc}{subsection}{1.1 The task statement}
The task asked for a new Mathematica notebook 7+1.nb in \texttt{D:\textbackslash{}Developer\textbackslash{}github\textbackslash{}7+1} defining a new realization of the octonions and a new representation of the Clifford algebra of an 8-dimensional vector space that mirror, as far as possible, the split octonions and the pseudo-Clifford algebra of 4+4.nb; real numbers and the real 8×8 and 16×16 matrices of 4+4.nb were to be used if possible, otherwise their complexification; examples were to check all work; notation was to be changed and documented where confusing. It fixed the coordinates
\begin{align*}
\text{4+4.nb:}\quad &X=(x_0,\dots,x_7),\ x_0=\text{hidden space-like},\ x_1,x_2,x_3=\text{3-space}, \\
&\phantom{X=(x_0,\dots,x_7),\ }x_4=\text{time},\ x_5,x_6,x_7=\text{auxiliary time-like}, \\
\text{7+1.nb:}\quad &X=(x_0,\dots,x_7),\ x_0=\text{time (octonion unit)},\ x_1,x_2,x_3=\text{3-space}, \\
&\phantom{X=(x_0,\dots,x_7),\ }x_4,\dots,x_7=\text{auxiliary space-like},
\end{align*}
and gave two definitions. [1A] the octonions: $e_0=1$, $e_i^2=-1$, Fano plane, Cayley–Dickson doubling $\mathbb O=\mathbb H\oplus\mathbb H$, conjugation $\bar x=x_0-\sum_ix_ie_i$ with $\overline{xy}=\bar y\bar x$, $N(x)=x\bar x=\sum_Ax_A^2$, $N(xy)=N(x)N(y)$, $x^{-1}=\bar x/N(x)$, alternativity, power-associativity. [1B] the Clifford algebra $Cl(V,Q)=\mathcal T(V)/\mathcal I_Q(V)$ with $v^2=Q(v)1$, $vw+wv=2B(v,w)1$, $B(v,w)=\frac12(Q(v+w)-Q(v)-Q(w))$, dimension $2^8=256$, a classification table for $p+q=8$, and $Cl_8(\mathbb C)\cong\mathbb C(16)$. Section 4 lists the statements of [1A] and [1B] that had to be corrected.

\Needspace{8\baselineskip}\subsection*{1.2 The notebook 4+4.nb}\addcontentsline{toc}{subsection}{1.2 The notebook 4+4.nb}
File \texttt{4+4.nb}, 9348656 bytes, written by Wolfram 15.0 (header \texttt{CreatedBy='Wolfram 15.0'}), SHA-256 \texttt{740023233eaa0cd24ca0e2e7010094671b1d314b698114c08db94050d3e0aed5}. It was read with the parser \texttt{tools/nbparse.py} (806 cells: 366 input, 366 output, 49 text). Every object used by 7+1.nb was reproduced with the formula of the cell given here, and each reproduction is checked in 7+1.nb section 2:

{\footnotesize
\begin{longtable}{>{\raggedright\arraybackslash}p{\dimexpr 0.1383\linewidth-2\tabcolsep\relax}>{\raggedright\arraybackslash}p{\dimexpr 0.2128\linewidth-2\tabcolsep\relax}>{\raggedright\arraybackslash}p{\dimexpr 0.6489\linewidth-2\tabcolsep\relax}}
\toprule
\textbf{object} & \textbf{4+4.nb cell} & \textbf{definition in 4+4.nb} \\
\midrule\endfirsthead
\toprule
\textbf{object} & \textbf{4+4.nb cell} & \textbf{definition in 4+4.nb} \\
\midrule\endhead
\endfoot
\bottomrule\endlastfoot
$\eta^{(4,4)}$ & In[9] & $\mathrm{diag}(1_4,-1_4)$ (\texttt{η4488}) \\
$\sigma$ & In[18] & $\left(\begin{smallmatrix}0&1_4\\1_4&0\end{smallmatrix}\right)$ (\texttt{σ=σ88}) \\
$s_h,t_h$ & In[29], In[35], In[36] & $(s_h)_{pq}=\epsilon_{hpq4}-(\delta_{p4}\delta_{qh}-\delta_{ph}\delta_{q4})$, $(t_h)_{pq}=\epsilon_{hpq4}+(\dots)$ \\
$\zeta_A$ & In[60]–In[63] & $\zeta_0=1_8$, $\zeta_{7-h}=\left(\begin{smallmatrix}0&t_h\\-t_h&0\end{smallmatrix}\right)$, $\zeta_h=\left(\begin{smallmatrix}0&s_h\\s_h&0\end{smallmatrix}\right)$, $\zeta_7=\zeta_1\cdots\zeta_6$ \\
relations of $\zeta$ & In[64], In[65], In[67] & $\sigma=\zeta_1\zeta_2\zeta_3$; $\frac12\lbrace\zeta_A,\zeta_B\rbrace=-\eta^{(4,4)}_{AB}1_8$; $-\eta^{(4,4)}_{AA}\zeta_A=\zeta_A^{\mathsf T}$ \\
$\tau_A$ & In[70] & $\tau_0=1_8$, $\tau_A=\zeta_A$ \\
$\bar\tau_A$ & In[74] & $\bar\tau_0=1_8$, $\bar\tau_A=\sigma\tau_A^{\mathsf T}\sigma$ (symbolized tau-bar) \\
relations of $\tau$ & In[77], In[87], In[89] & $\frac12(\tau_A\bar\tau_B+\tau_B\bar\tau_A)=\eta^{(4,4)}_{AB}1_8$; $\tau_1\cdots\tau_7=1_8$; $\sigma=\tau_4\tau_5\tau_6\tau_7$ \\
$T_A$ & In[98] & $T_A=\left(\begin{smallmatrix}0&\bar\tau_A\\\tau_A&0\end{smallmatrix}\right)$ (\texttt{T16\textasciicircum{}A[A1]}) \\
$T_8$, $\sigma_{16}$ & In[99], In[97], In[106] & $T_8=T_0\cdots T_7$; $\sigma_{16}=T_0T_1T_2T_3=\mathrm{diag}(-\sigma,\sigma)$ \\
Clifford relation & In[110] & $\frac12\lbrace T_A,T_B\rbrace=\eta^{(4,4)}_{AB}1_{16}$ \\
$P_L,P_R$ & In[112]–In[114] & $P_L=\frac12(1_{16}-T_8)$, $P_R=\frac12(1_{16}+T_8)$ \\
$U_h$ & In[257], In[259] & eigenvectors of $\sigma$ divided by $\sqrt2$ \\
$u$ & In[264], In[266] & \texttt{hUSE = 8}, $u=U_8=(1,0,0,0,1,0,0,0)/\sqrt2$ \\
$g$ & In[283] & $g_{AB}=\eta^{(4,4)}$ \\
$E_{44}$, $E_{44}^{-1}$ & In[294]–In[296] & $(E_{44})_{Ba}=(U_h^{\mathsf T}\sigma\bar\tau_B)_a$; $E_{44}^{-1}=(U_h^{\mathsf T}\sigma U_h)(\eta^{(4,4)}_{BB}\tau_BU_h)_B$ \\
$\varepsilon_B$ & In[313] & \texttt{EA = Array[eA, 8]} \\
split table & In[325], Out[325] & $\varepsilon_B\varepsilon_C=\sum_D(g_{BB}E_{44}\tau_{B-1}E_{44}^{-1})_{DC}\varepsilon_D$ (\texttt{gridx}) \\*
\end{longtable}}
\Needspace{8\baselineskip}\subsection*{1.3 Literature}\addcontentsline{toc}{subsection}{1.3 Literature}
\begin{itemize}
\item J. C. Baez, \emph{The Octonions}, Bull. Amer. Math. Soc. 39 (2002) 145–205, arXiv:math/0105155, section 2.2 (the page https://math.ucr.edu/home/baez/octonions/node5.html, which the task statement itself cites, was read on 2026-10-26): the Cayley–Dickson product $(a,b)(c,d)=(ac-db^{*},\ a^{*}d+cb)$, where Baez writes $x^{*}$ for the conjugate $\bar x$ and the conjugate of a pair is $(a,b)^{*}=(a^{*},-b)$. 7+1.nb proves that its octonions satisfy this product with $x=a+\ell b$, $\ell=e_4$.
\item H. B. Lawson and M.-L. Michelsohn, \emph{Spin Geometry}, Princeton (1989): classification of real Clifford algebras, $Cl(p+1,q+1)\cong Cl(p,q)\otimes\mathbb R(2)$. 7+1.nb does not rely on it: it computes the classification directly (section 2.10 below).
\item Standard facts used in proofs: Schur's lemma (the intertwiners of an irreducible module form a division algebra; over $\mathbb C$, the scalars); $Cl_8(\mathbb C)\cong\mathbb C(16)$ has a unique irreducible module $\mathbb C^{16}$; Hurwitz's theorem (a unital composition algebra of dimension 8 is the octonions if its norm is positive definite and the split octonions if the norm is indefinite); Degen's eight-square identity.
\end{itemize}
\Needspace{12\baselineskip}\section*{2. Derivations: how each new object was obtained}\addcontentsline{toc}{section}{2. Derivations: how each new object was obtained}
\Needspace{8\baselineskip}\subsection*{2.1 Octonion matrices \texorpdfstring{$\tau^{O}$}{τO}}\addcontentsline{toc}{subsection}{2.1 Octonion matrices \texorpdfstring{$\tau^{O}$}{τO}}
Origin: the split matrices of 4+4.nb and the observation (checked) that $\sigma=\tau_1\tau_2\tau_3=\tau_4\tau_5\tau_6\tau_7$ commutes with $\tau_1,\tau_2,\tau_3$, anticommutes with $\tau_4,\dots,\tau_7$ and squares to $1_8$. Hence
\begin{align*}
\tau^{O}_0&=1_8,\quad \tau^{O}_A=\tau_A\ (A\le3),\quad \tau^{O}_k=\tau_k\sigma\ (k\ge4),\qquad (\tau_k\sigma)^2=\tau_k\sigma\tau_k\sigma=-\tau_k^2\sigma^2=-1_8, \\
\tau^{O}_h\tau^{O}_k&=\tau_h\tau_k\sigma=-\tau_k\tau_h\sigma=-\tau_k\sigma\tau_h=-\tau^{O}_k\tau^{O}_h,\qquad \tau^{O}_k\tau^{O}_l=\tau_k\sigma\tau_l\sigma=-\tau_k\tau_l=\tau_l\tau_k=-\tau^{O}_l\tau^{O}_k\quad(k\neq l\ge4).
\end{align*}
This is the real counterpart of the complex substitution $\tau_k\to i\tau_k$ that turns split octonions into octonions over $\mathbb C$.

\Needspace{8\baselineskip}\subsection*{2.2 Octonion vielbein and multiplication constants}\addcontentsline{toc}{subsection}{2.2 Octonion vielbein and multiplication constants}
Origin: In[294] and In[325] of 4+4.nb with the split spinor metric $\sigma$ replaced by the Euclidean $g^{O}=1_8$ and the same unit spinor $u$:
\begin{align*}
(E^{O})_{Ba}&=\big(u^{\mathsf T}g^{O}\bar\tau^{O}_B\big)_a,\qquad L^{O}_A=g^{O}_{AA}E^{O}\tau^{O}_A(E^{O})^{-1},\qquad e_Ae_K=\sum_L(L^{O}_A)_{LK}e_L .
\end{align*}
Why the result is the octonion algebra: $(E^{O})^{-1}e_0=u$ makes $e_0$ a two-sided unit, and $\tau^{O}(x)\bar\tau^{O}(x)=N(x)1_8$ together with the orthogonality of $E^{O}$ gives $N(xy)=N(x)N(y)$. A unital composition algebra of dimension 8 with a positive definite norm is the octonions (Hurwitz). The notebook checks all the laws directly.

\Needspace{8\baselineskip}\subsection*{2.3 Cayley–Dickson form}\addcontentsline{toc}{subsection}{2.3 Cayley–Dickson form}
Origin: Baez 2002, section 2.2. With $\ell=e_4$, $x=a+\ell b$, $a=(x_0,\dots,x_3)$, $b=(x_4,\dots,x_7)$:
\begin{equation*}
(a,b)(c,d)=\big(ac+\gamma\,d\bar b,\ \bar a d+cb\big),\qquad \gamma=-1\ (\text{7+1.nb}),\qquad \gamma=+1\ (\text{4+4.nb}).
\end{equation*}
\Needspace{8\baselineskip}\subsection*{2.4 Euclidean mirror \texorpdfstring{$T^{O}$}{TO}}\addcontentsline{toc}{subsection}{2.4 Euclidean mirror \texorpdfstring{$T^{O}$}{TO}}
Origin: In[98] of 4+4.nb with $\tau\to\tau^{O}$: $T^{O}_A=\left(\begin{smallmatrix}0&\bar\tau^{O}_A\\\tau^{O}_A&0\end{smallmatrix}\right)$, so that $\lbrace T^{O}_A,T^{O}_B\rbrace=\left(\begin{smallmatrix}\bar\tau^{O}_A\tau^{O}_B+\bar\tau^{O}_B\tau^{O}_A&0\\0&\tau^{O}_A\bar\tau^{O}_B+\tau^{O}_B\bar\tau^{O}_A\end{smallmatrix}\right)=2\delta_{AB}1_{16}$. In terms of 4+4.nb: $T^{O}_A=T_A$ ($A\le3$), $T^{O}_k=T_kS_{44}$ ($k\ge4$), $S_{44}=T_4T_5T_6T_7=\mathrm{diag}(\sigma,\sigma)$.

\Needspace{8\baselineskip}\subsection*{2.5 Real 7+1 generators \texorpdfstring{$\Gamma^{(1,7)}$}{Γ(1,7)}}\addcontentsline{toc}{subsection}{2.5 Real 7+1 generators \texorpdfstring{$\Gamma^{(1,7)}$}{Γ(1,7)}}
Origin: the standard construction of a Lorentzian Clifford algebra from a Euclidean one, applied to $T^{O}$:
\begin{align*}
\Gamma^{(1,7)}_0&=T^{O}_0,\qquad \Gamma^{(1,7)}_i=T^{O}_0T^{O}_i,\qquad (T^{O}_0T^{O}_i)^2=-(T^{O}_0)^2(T^{O}_i)^2=-1_{16}, \\
\lbrace T^{O}_0,T^{O}_0T^{O}_i\rbrace&=T^{O}_i+T^{O}_0T^{O}_iT^{O}_0=T^{O}_i-T^{O}_i=0,\qquad \lbrace T^{O}_0T^{O}_i,T^{O}_0T^{O}_j\rbrace=-\lbrace T^{O}_i,T^{O}_j\rbrace=-2\delta_{ij}1_{16}.
\end{align*}
Hence $\lbrace\Gamma^{(1,7)}_A,\Gamma^{(1,7)}_B\rbrace=2\eta^{(1,7)}_{AB}1_{16}$, and $\Gamma^{(1,7)}(x)^2=\mathrm{Re}(xx)1_{16}$.

\Needspace{8\baselineskip}\subsection*{2.6 Complexified generators \texorpdfstring{$\Gamma^{(7,1)}$}{Γ(7,1)}}\addcontentsline{toc}{subsection}{2.6 Complexified generators \texorpdfstring{$\Gamma^{(7,1)}$}{Γ(7,1)}}
Origin: the task's instruction to use the complexification of the 4+4 matrices if real numbers fail. Only the time block is multiplied by $i$: $\tau^{C}_0=\bar\tau^{C}_0=i1_8$, $\tau^{C}_k=\tau^{O}_k$; then $\Gamma^{(7,1)}_A=\left(\begin{smallmatrix}0&\bar\tau^{C}_A\\\tau^{C}_A&0\end{smallmatrix}\right)$, $(\Gamma^{(7,1)}_0)^2=i^2\,1_{16}=-1_{16}$, and $\Gamma^{(7,1)}_A=S(i\Gamma^{(1,7)}_A)S^{-1}$ with $S=(1_{16}-i\Gamma^{(1,7)}_0)/\sqrt2$.

\Needspace{8\baselineskip}\subsection*{2.7 The impossibility theorem}\addcontentsline{toc}{subsection}{2.7 The impossibility theorem}
Origin: Schur's lemma and the uniqueness of the irreducible module of $Cl_8(\mathbb C)\cong\mathbb C(16)$. The intertwiners $M$ with $M\overline{\Gamma^{(7,1)}_A}=\Gamma^{(7,1)}_AM$ are computed in 7+1.nb with \texttt{NullSpace}; they form the line through $M_0=\Gamma^{(1,7)}_8=\left(\begin{smallmatrix}0&-1_8\\1_8&0\end{smallmatrix}\right)$, and $M_0\overline{M_0}=-1_{16}$. A real realization $R_A=T\Gamma^{(7,1)}_AT^{-1}$ would give an intertwiner $M=T^{-1}\overline T$ with $M\overline M=1_{16}$, impossible because $M=cM_0$ gives $M\overline M=-|c|^21_{16}$.

\Needspace{8\baselineskip}\subsection*{2.8 Real 32×32 generators}\addcontentsline{toc}{subsection}{2.8 Real 32×32 generators}
Origin: realification. $\widehat\Gamma^{(7,1)}_A=\varepsilon\otimes\Gamma^{(1,7)}_A$ with $\varepsilon^2=-1_2$ gives $\lbrace\widehat\Gamma^{(7,1)}_A,\widehat\Gamma^{(7,1)}_B\rbrace=\varepsilon^2\otimes2\eta^{(1,7)}_{AB}1_{16}=2\eta^{(7,1)}_{AB}1_{32}$. Commutant: $Q_i=\varepsilon\otimes1_{16}$, $Q_j=\sigma_1\otimes\Gamma^{(1,7)}_8$, $Q_k=\sigma_3\otimes\Gamma^{(1,7)}_8$, found from $\varepsilon\sigma_1=-\sigma_1\varepsilon$, $\varepsilon\sigma_3=-\sigma_3\varepsilon$ and $\Gamma^{(1,7)}_8\Gamma^{(1,7)}_A=-\Gamma^{(1,7)}_A\Gamma^{(1,7)}_8$.

\Needspace{8\baselineskip}\subsection*{2.9 Examples}\addcontentsline{toc}{subsection}{2.9 Examples}
so(1,7): $\Sigma_{AB}=\frac14[\Gamma^{(1,7)}_A,\Gamma^{(1,7)}_B]$; the commutation relations follow from the Clifford relation. Boost and rotation: $(\Gamma^{(1,7)}_0\Gamma^{(1,7)}_1)^2=1_{16}$, $(\Gamma^{(1,7)}_4\Gamma^{(1,7)}_5)^2=-1_{16}$, so $e^{\beta\Sigma_{01}}=\cosh\frac\beta2+\sinh\frac\beta2\Gamma^{(1,7)}_0\Gamma^{(1,7)}_1$ and $e^{\theta\Sigma_{45}}=\cos\frac\theta2+\sin\frac\theta2\Gamma^{(1,7)}_4\Gamma^{(1,7)}_5$. Dirac operator: $(\Gamma^A\partial_A)^2=\frac12\lbrace\Gamma^A,\Gamma^B\rbrace\partial_A\partial_B=\eta_{(1,7)}^{AB}\partial_A\partial_B$. Real amplitudes: $p=(3,2,2,0,0,0,0,0)$, $p^{\mathsf T}\eta^{(1,7)}p=1$, $\not p^{\,2}=1_{16}$; the real matrix $\not p$ has a real 8-dimensional eigenspace for the eigenvalue $+1$.

Real spinor fields. If $J$ anticommutes with every $\Gamma^{(1,7)}_A$, then $J\Gamma^{(1,7)}_8$ commutes with every $\Gamma^{(1,7)}_A$ (because $\Gamma^{(1,7)}_8$ anticommutes with each of them), so by irreducibility $J\Gamma^{(1,7)}_8=c\,1_{16}$ and $J=-c\,\Gamma^{(1,7)}_8$ (as $(\Gamma^{(1,7)}_8)^{-1}=-\Gamma^{(1,7)}_8$), whose square is $-c^2\,1_{16}$. Hence the real operator $\Gamma^A\partial_A-mJ$ has the square $\partial_0^2-\dots-\partial_7^2-c^2m^2$ (a tachyon), and a real equation $(G^A\partial_A-m)\Psi=0$ implies the Klein–Gordon equation only if the real $G_A$ obey the $\eta^{(7,1)}$ relation, impossible for 16×16 matrices (2.7). Massless: for $k=(1,1,0,\dots,0)$, $\not k^{\,2}=k^{\mathsf T}\eta^{(1,7)}k\,1_{16}=0$ and $\Gamma^A\partial_A[w\cos(k\cdot x)]=-\sin(k\cdot x)\,\not k\,w$, which vanishes for the 8 real solutions of $\not k\,w=0$. Massive, with the real 32×32 matrices: $P=\sum_Ap^A\widehat\Gamma^{(7,1)}_A$, $P^2=-m^2\,1_{32}$, and with $\phi=p\cdot x$, $\widehat\Gamma^A=\eta_{(7,1)}^{AA}\widehat\Gamma^{(7,1)}_A$:
\begin{align*}
\widehat\Gamma^{A}\partial_A\Big(w\cos\phi+\tfrac1mPw\sin\phi\Big)&=-P\Big(-w\sin\phi+\tfrac1mPw\cos\phi\Big) \\
&=Pw\sin\phi+m\,w\cos\phi=m\Big(w\cos\phi+\tfrac1mPw\sin\phi\Big),
\end{align*}
using $\partial_A\phi=(\eta^{(1,7)}p)_A$ and $\eta_{(7,1)}^{AA}\eta^{(1,7)}_{AA}=-1$.

\Needspace{8\baselineskip}\subsection*{2.10 Classification of the nine signatures}\addcontentsline{toc}{subsection}{2.10 Classification of the nine signatures}
For each $p$, complex generators $G_A=T^{O}_A$ ($A<p$), $G_A=iT^{O}_A$ ($A\ge p$) satisfy $\lbrace G_A,G_B\rbrace=\pm2\delta_{AB}1_{16}$ with $p$ plus signs. The indicator sign of $M\overline M$ ($M\overline{G_A}=G_AM$) is $+1$ for $p\in\lbrace8,5,4,1,0\rbrace$ ($\mathbb R(16)$) and $-1$ for $p\in\lbrace7,6,3,2\rbrace$ ($\mathbb H(8)$). An independent numerical method (\texttt{tools/classify\_clifford.py}: commutant of the realified 32-dimensional representation, whose traceless part has a negative definite trace form exactly in the quaternionic case) gives the same table. Explicit real generators: $T^{O}$ (8,0), $\Gamma^{(5,3)}_0=T_0$, $\Gamma^{(5,3)}_i=T_0T_i$ (5,3) with $\eta^{(5,3)}=\mathrm{diag}(1,-1,-1,-1,1,1,1,1)$, $T$ (4,4), $\Gamma^{(1,7)}$ (1,7), $\Gamma^{(0,8)}_A=T^{O}_8T^{O}_A$ (0,8).

\Needspace{12\baselineskip}\section*{3. Corrections to the task statement (with the evidence in 7+1.nb)}\addcontentsline{toc}{section}{3. Corrections to the task statement (with the evidence in 7+1.nb)}
{\footnotesize
\begin{longtable}{>{\raggedright\arraybackslash}p{\dimexpr 0.3191\linewidth-2\tabcolsep\relax}>{\raggedright\arraybackslash}p{\dimexpr 0.3617\linewidth-2\tabcolsep\relax}>{\raggedright\arraybackslash}p{\dimexpr 0.3191\linewidth-2\tabcolsep\relax}}
\toprule
\textbf{statement in the task} & \textbf{correction} & \textbf{evidence (check in 7+1.nb)} \\
\midrule\endfirsthead
\toprule
\textbf{statement in the task} & \textbf{correction} & \textbf{evidence (check in 7+1.nb)} \\
\midrule\endhead
\endfoot
\bottomrule\endlastfoot
{}[1A]: $(a,b)(c,d)=(ac-\bar d b,\ \bar a d+cb)$ defines the octonions & not alternative, norm not multiplicative; correct form $(ac-d\bar b,\ \bar a d+cb)$ (Baez 2002, 2.2) & "the formula ... of [1A] is NOT alternative", "... its norm is not multiplicative", "Cayley-Dickson, \texorpdfstring{\ensuremath{\gamma}}{γ} = -1 ..." \\
{}[1A]: first formula $(ac-\bar d b,\ da+cb\bar x)$ & contains a stray "x"; not interpretable & — \\
{}[1B]: $1+8+28+56+72+56+28+8+1=256$ & $\binom84=70$ & "grades: binomial(8,k) == ..." \\
{}[1B]: $Cl(7,1)\cong\mathbb R(16)$ & $Cl(7,1)\cong\mathbb H(8)$ for $v^2=Q(v)$ & "indicator(\texorpdfstring{\ensuremath{\Gamma}}{Γ}71) == -1 ...", "classification: ..." \\
{}[1B]: $Cl(5,3)\cong\mathbb C(16)$, $Cl(3,5)\cong\mathbb C(16)$ & $Cl(5,3)\cong\mathbb R(16)$, $Cl(3,5)\cong\mathbb H(8)$; $\mathbb C(16)$ has real dimension 512 & "classification: ...", "\texorpdfstring{\ensuremath{\mathbb{C}}}{ℂ}(16) cannot occur ..." \\
{}[1B]: metric-sign strings for (6,2), (5,3), (4,4), (2,6) & mistyped (wrong number of signs) & table of section 5.6 of 7+1.nb \\
split octonions are a "pseudo-normed division algebra" & they have zero divisors: $(\varepsilon_1+\varepsilon_5)(\varepsilon_1-\varepsilon_5)=0$ & "split: (eA[1] + eA[5])(eA[1] - eA[5]) == 0 ..." \\*
\end{longtable}}
\Needspace{12\baselineskip}\section*{4. Verification record}\addcontentsline{toc}{section}{4. Verification record}
\Needspace{8\baselineskip}\subsection*{4.1 The Mathematica run}\addcontentsline{toc}{subsection}{4.1 The Mathematica run}
{\footnotesize
\begin{longtable}{>{\raggedright\arraybackslash}p{\dimexpr 0.3617\linewidth-2\tabcolsep\relax}>{\raggedright\arraybackslash}p{\dimexpr 0.6383\linewidth-2\tabcolsep\relax}}
\toprule
\textbf{item} & \textbf{value} \\
\midrule\endfirsthead
\toprule
\textbf{item} & \textbf{value} \\
\midrule\endhead
\endfoot
\bottomrule\endlastfoot
program & \texttt{C:\textbackslash{}Program Files\textbackslash{}Wolfram Research\textbackslash{}Wolfram\textbackslash{}15.0.1\textbackslash{}WolframNB.exe} \\
\$Version & 15.0.1 for Microsoft Windows (64-bit) (July 2, 2026) \\
method & \texttt{Get["D:/Developer/github/7+1/tools/run\_eval.wl"]}: \texttt{NotebookEvaluate[nb, InsertResults -> True]}, then \texttt{NotebookSave} \\
finished & Mon 26 Oct 2026 14:54:43 \\
evaluation time & 5.09545 s \\
private context of the notebook & \texttt{Notebook\$\$43\$40478} \\
input cells / output cells / message cells & 56 / 154 / 0 \\
checks / True & 128 / 128 \\
commutant computation (8192×1024 sparse \texttt{NullSpace}) & 1.0349715 s \\
SHA-256 of the evaluated 7+1.nb (306963 bytes) & \texttt{f7abe7279b5540e914de876dc055c199bd491b9c2d9522bf28dbc23689e237ce} \\*
\end{longtable}}
\Needspace{8\baselineskip}\subsection*{4.2 All checks of 7+1.nb and their results}\addcontentsline{toc}{subsection}{4.2 All checks of 7+1.nb and their results}
Each row is one cell \texttt{check["name", test]} of the notebook, in order, grouped by the section of the notebook in which it stands; the test is the Wolfram Language expression in that cell, and the result is what Wolfram 15.0.1 returned.

\Needspace{6\baselineskip}\subsubsection*{2.1 \texorpdfstring{\ensuremath{\sigma}}{σ}, the self-dual and anti-self-dual blocks, \texorpdfstring{\ensuremath{\zeta}}{ζ}, \texorpdfstring{\ensuremath{\tau}}{τ} and \texorpdfstring{\ensuremath{\tau}}{τ}bar}
{\footnotesize
\begin{longtable}{>{\raggedright\arraybackslash}p{\dimexpr 0.0700\linewidth-2\tabcolsep\relax}>{\raggedright\arraybackslash}p{\dimexpr 0.8100\linewidth-2\tabcolsep\relax}>{\raggedright\arraybackslash}p{\dimexpr 0.1200\linewidth-2\tabcolsep\relax}}
\toprule
\textbf{\#} & \textbf{check} & \textbf{result} \\
\midrule\endfirsthead
\toprule
\textbf{\#} & \textbf{check} & \textbf{result} \\
\midrule\endhead
\endfoot
\bottomrule\endlastfoot
1 & \texttt{4+4: ζ[7] == diag(-1,-1,-1,-1,1,1,1,1)} & True \\
2 & \texttt{4+4: σ == ζ[1].ζ[2].ζ[3]} & True \\
3 & \texttt{4+4: (ζ[A].ζ[B] + ζ[B].ζ[A])/2 == -η4488[[A+1,B+1]] ID8 (A,B = 1..7)} & True \\
4 & \texttt{4+4: -η4488[[A+1,A+1]] ζ[A] == Transpose[ζ[A]] (A = 1..7)} & True \\
5 & \texttt{4+4: (τ[A].τbar[B] + τ[B].τbar[A])/2 == η4488[[A+1,B+1]] ID8} & True \\
6 & \texttt{4+4: (τbar[A].τ[B] + τbar[B].τ[A])/2 == η4488[[A+1,B+1]] ID8} & True \\
7 & \texttt{4+4: τ[1].τ[2]...τ[7] == ID8} & True \\
8 & \texttt{4+4: σ == τ[1].τ[2].τ[3] == τ[4].τ[5].τ[6].τ[7]} & True \\
9 & \texttt{4+4: σ commutes with τ[1..3] and anticommutes with τ[4..7]; σ.σ == ID8} & True \\*
\end{longtable}}
\Needspace{6\baselineskip}\subsubsection*{2.2 The 4+4 Dirac matrices T16 and the chiral projectors}
{\footnotesize
\begin{longtable}{>{\raggedright\arraybackslash}p{\dimexpr 0.0700\linewidth-2\tabcolsep\relax}>{\raggedright\arraybackslash}p{\dimexpr 0.8100\linewidth-2\tabcolsep\relax}>{\raggedright\arraybackslash}p{\dimexpr 0.1200\linewidth-2\tabcolsep\relax}}
\toprule
\textbf{\#} & \textbf{check} & \textbf{result} \\
\midrule\endfirsthead
\toprule
\textbf{\#} & \textbf{check} & \textbf{result} \\
\midrule\endhead
\endfoot
\bottomrule\endlastfoot
10 & \texttt{4+4: (T16[A].T16[B] + T16[B].T16[A])/2 == η4488[[A+1,B+1]] ID16} & True \\
11 & \texttt{4+4: T16[A] real (entries 0, ±1)} & True \\
12 & \texttt{4+4: T16[8] == diag(-ID8, ID8)} & True \\
13 & \texttt{4+4: σ16 == diag(-σ, σ)} & True \\
14 & \texttt{4+4: PL44 + PR44 == ID16, PL44.PL44 == PL44, PR44.PR44 == PR44, PL44.PR44 == 0} & True \\*
\end{longtable}}
\Needspace{6\baselineskip}\subsubsection*{2.3 The split-octonion multiplication constants (gridx of 4+4.nb) and how they come from \texorpdfstring{\ensuremath{\tau}}{τ}}
{\footnotesize
\begin{longtable}{>{\raggedright\arraybackslash}p{\dimexpr 0.0700\linewidth-2\tabcolsep\relax}>{\raggedright\arraybackslash}p{\dimexpr 0.8100\linewidth-2\tabcolsep\relax}>{\raggedright\arraybackslash}p{\dimexpr 0.1200\linewidth-2\tabcolsep\relax}}
\toprule
\textbf{\#} & \textbf{check} & \textbf{result} \\
\midrule\endfirsthead
\toprule
\textbf{\#} & \textbf{check} & \textbf{result} \\
\midrule\endhead
\endfoot
\bottomrule\endlastfoot
15 & \texttt{4+4: uBasis == evecs/Sqrt[2] from Eigensystem[σ] (the basis of 4+4.nb)} & True \\
16 & \texttt{4+4: unit == (1,0,0,0,1,0,0,0)/Sqrt[2], σ.unit == unit, unit.unit == 1, unit.σ.unit == 1} & True \\
17 & \texttt{4+4 (In[294]): E44inv[h] == (u.σ.u) Transpose[Table[η4488[[B,B]] τ[B-1].u, \{B,1,8\}]] for h = 1..8} & True \\
18 & \texttt{4+4: gridx44 is identical to the split-octonion table printed by 4+4.nb} & True \\
19 & \texttt{4+4: LS[A] are signed permutation matrices (constants 0, ±1)} & True \\*
\end{longtable}}
\Needspace{6\baselineskip}\subsubsection*{2.4 Split octonions: composition algebra, but not a division algebra}
{\footnotesize
\begin{longtable}{>{\raggedright\arraybackslash}p{\dimexpr 0.0700\linewidth-2\tabcolsep\relax}>{\raggedright\arraybackslash}p{\dimexpr 0.8100\linewidth-2\tabcolsep\relax}>{\raggedright\arraybackslash}p{\dimexpr 0.1200\linewidth-2\tabcolsep\relax}}
\toprule
\textbf{\#} & \textbf{check} & \textbf{result} \\
\midrule\endfirsthead
\toprule
\textbf{\#} & \textbf{check} & \textbf{result} \\
\midrule\endhead
\endfoot
\bottomrule\endlastfoot
20 & \texttt{split: eA[1] is the identity; eA[2..4]\textasciicircum{}2 == -eA[1]; eA[5..8]\textasciicircum{}2 == +eA[1]} & True \\
21 & \texttt{split: x x̄ == N(x) eA[1] with N = η4488 (symbolic)} & True \\
22 & \texttt{split: N(xy) == N(x) N(y) (symbolic)} & True \\
23 & \texttt{split: alternative x(xy) == (xx)y and (yx)x == y(xx) (symbolic)} & True \\
24 & \texttt{split: (eA[1] + eA[5])(eA[1] - eA[5]) == 0 with both factors nonzero: NOT a division algebra} & True \\*
\end{longtable}}
\Needspace{6\baselineskip}\subsubsection*{3.1 From the split matrices \texorpdfstring{\ensuremath{\tau}}{τ} to the octonionic matrices \texorpdfstring{\ensuremath{\tau}}{τ}O}
{\footnotesize
\begin{longtable}{>{\raggedright\arraybackslash}p{\dimexpr 0.0700\linewidth-2\tabcolsep\relax}>{\raggedright\arraybackslash}p{\dimexpr 0.8100\linewidth-2\tabcolsep\relax}>{\raggedright\arraybackslash}p{\dimexpr 0.1200\linewidth-2\tabcolsep\relax}}
\toprule
\textbf{\#} & \textbf{check} & \textbf{result} \\
\midrule\endfirsthead
\toprule
\textbf{\#} & \textbf{check} & \textbf{result} \\
\midrule\endhead
\endfoot
\bottomrule\endlastfoot
25 & \texttt{τO[A] (A = 1..7): real, antisymmetric, τO[A].τO[A] == -ID8} & True \\
26 & \texttt{τO[A] anticommute: τO[A].τO[B] == -τO[B].τO[A] (A != B, A,B = 1..7)} & True \\
27 & \texttt{mirror of 4+4: (τO[A].τObar[B] + τO[B].τObar[A])/2 == gO[[A+1,B+1]] ID8} & True \\
28 & \texttt{mirror of 4+4: (τObar[A].τO[B] + τObar[B].τO[A])/2 == gO[[A+1,B+1]] ID8} & True \\
29 & \texttt{mirror of 4+4: τO[1].τO[2]...τO[7] == ID8} & True \\
30 & \texttt{mirror of 4+4: τObar[A] == gO.Transpose[τO[A]].gO, as τbar[A] == σ.Transpose[τ[A]].σ} & True \\
31 & \texttt{τO[4..7] are products of three split τ: τO[4] == τ[5].τ[6].τ[7], τO[5] == -τ[4].τ[6].τ[7], τO[6] == τ[4].τ[5].τ[7], τO[7] == -τ[4].τ[5].τ[6]} & True \\*
\end{longtable}}
\Needspace{6\baselineskip}\subsubsection*{3.2 The octonion vielbein (the same unit spinor u as in 4+4.nb)}
{\footnotesize
\begin{longtable}{>{\raggedright\arraybackslash}p{\dimexpr 0.0700\linewidth-2\tabcolsep\relax}>{\raggedright\arraybackslash}p{\dimexpr 0.8100\linewidth-2\tabcolsep\relax}>{\raggedright\arraybackslash}p{\dimexpr 0.1200\linewidth-2\tabcolsep\relax}}
\toprule
\textbf{\#} & \textbf{check} & \textbf{result} \\
\midrule\endfirsthead
\toprule
\textbf{\#} & \textbf{check} & \textbf{result} \\
\midrule\endhead
\endfoot
\bottomrule\endlastfoot
32 & \texttt{σ.unit == unit, hence τO[A].unit == τ[A].unit (A = 0..7)} & True \\
33 & \texttt{EO is orthogonal: EOinv == Transpose[EO]} & True \\
34 & \texttt{mirror of 4+4.nb In[294]: EOinv == (u.gO.u) Transpose[Table[gO[[B,B]] τO[B-1].u, \{B,1,8\}]]} & True \\
35 & \texttt{EO == η4488.E44[hUSE]: the two vielbeins differ only by the signs of rows 5..8} & True \\*
\end{longtable}}
\Needspace{6\baselineskip}\subsubsection*{3.3 The octonion multiplication constants: gridO, the mirror of gridx}
{\footnotesize
\begin{longtable}{>{\raggedright\arraybackslash}p{\dimexpr 0.0700\linewidth-2\tabcolsep\relax}>{\raggedright\arraybackslash}p{\dimexpr 0.8100\linewidth-2\tabcolsep\relax}>{\raggedright\arraybackslash}p{\dimexpr 0.1200\linewidth-2\tabcolsep\relax}}
\toprule
\textbf{\#} & \textbf{check} & \textbf{result} \\
\midrule\endfirsthead
\toprule
\textbf{\#} & \textbf{check} & \textbf{result} \\
\midrule\endhead
\endfoot
\bottomrule\endlastfoot
36 & \texttt{gridO equals the octonion table computed independently in Python} & True \\
37 & \texttt{LO[A] are signed permutation matrices (all multiplication constants are 0, ±1)} & True \\
38 & \texttt{oMul reproduces gridO on all 64 pairs of basis units} & True \\*
\end{longtable}}
\Needspace{6\baselineskip}\subsubsection*{3.4 The properties listed in [1A], verified}
{\footnotesize
\begin{longtable}{>{\raggedright\arraybackslash}p{\dimexpr 0.0700\linewidth-2\tabcolsep\relax}>{\raggedright\arraybackslash}p{\dimexpr 0.8100\linewidth-2\tabcolsep\relax}>{\raggedright\arraybackslash}p{\dimexpr 0.1200\linewidth-2\tabcolsep\relax}}
\toprule
\textbf{\#} & \textbf{check} & \textbf{result} \\
\midrule\endfirsthead
\toprule
\textbf{\#} & \textbf{check} & \textbf{result} \\
\midrule\endhead
\endfoot
\bottomrule\endlastfoot
39 & \texttt{e[0] is a two-sided identity} & True \\
40 & \texttt{e[i] e[i] == -e[0] (i = 1..7)} & True \\
41 & \texttt{e[i] e[j] == -e[j] e[i] (i != j, i,j = 1..7): non-commutative} & True \\
42 & \texttt{7 Fano lines; every pair of distinct imaginary units lies on exactly one line} & True \\
43 & \texttt{each oriented line (a,b,c) is cyclic: e[a]e[b] == e[c], e[b]e[c] == e[a], e[c]e[a] == e[b]} & True \\
44 & \texttt{each line spans a quaternion subalgebra: \{e[0], e[a], e[b], e[c]\} closed and associative} & True \\
45 & \texttt{the picture's 7 lines are exactly fanoLines} & True \\
46 & \texttt{left alternative: x(xy) == (xx)y (symbolic)} & True \\
47 & \texttt{right alternative: (yx)x == y(xx) (symbolic)} & True \\
48 & \texttt{flexible: (xy)x == x(yx) (symbolic)} & True \\
49 & \texttt{associator is alternating: [x,y,z] == -[y,x,z] == -[x,z,y] (symbolic)} & True \\
50 & \texttt{power associativity: x(xx) == (xx)x (symbolic)} & True \\
51 & \texttt{Moufang: z(x(zy)) == ((zx)z)y, x(z(yz)) == ((xz)y)z, (zx)(yz) == (z(xy))z (x, y basis units, z symbolic)} & True \\
52 & \texttt{x x̄ == x̄ x == N(x) e[0] (symbolic)} & True \\
53 & \texttt{conjugation is an anti-automorphism: conj(xy) == conj(y) conj(x) (symbolic)} & True \\
54 & \texttt{Re(x) == (x + x̄)/2 == x0} & True \\
55 & \texttt{norm is multiplicative: N(xy) == N(x) N(y) (symbolic; Degen's eight-square identity)} & True \\
56 & \texttt{inverse x\textasciicircum{}-1 = x̄/N(x): x x\textasciicircum{}-1 == x\textasciicircum{}-1 x == e[0] (symbolic)} & True \\
57 & \texttt{every x satisfies x x - 2 Re(x) x + N(x) e[0] == 0, so the subalgebra generated by x is ℝ or ℂ} & True \\
58 & \texttt{exactly 168 of the 512 ordered basis triples are non-associative (= 7*6*4: distinct imaginary units not on one Fano line)} & True \\
59 & \texttt{example: (e[1]e[2])e[4] == e[7] but e[1](e[2]e[4]) == -e[7]} & True \\
60 & \texttt{numerical example: N(x y) == N(x) N(y) == 15*6 == 90 and x x\textasciicircum{}-1 == e[0]} & True \\
61 & \texttt{numerical example: x y != y x and (x y) z != x (y z)} & True \\*
\end{longtable}}
\Needspace{6\baselineskip}\subsubsection*{3.5 Cayley–Dickson doubling: \texorpdfstring{\ensuremath{\mathbb{O}}}{𝕆} = \texorpdfstring{\ensuremath{\mathbb{H}}}{ℍ} \texorpdfstring{\ensuremath{\oplus}}{⊕} \texorpdfstring{\ensuremath{\ell}}{ℓ}\texorpdfstring{\ensuremath{\mathbb{H}}}{ℍ}}
{\footnotesize
\begin{longtable}{>{\raggedright\arraybackslash}p{\dimexpr 0.0700\linewidth-2\tabcolsep\relax}>{\raggedright\arraybackslash}p{\dimexpr 0.8100\linewidth-2\tabcolsep\relax}>{\raggedright\arraybackslash}p{\dimexpr 0.1200\linewidth-2\tabcolsep\relax}}
\toprule
\textbf{\#} & \textbf{check} & \textbf{result} \\
\midrule\endfirsthead
\toprule
\textbf{\#} & \textbf{check} & \textbf{result} \\
\midrule\endhead
\endfoot
\bottomrule\endlastfoot
62 & \texttt{ℍ = span\{e[0],e[1],e[2],e[3]\} is a subalgebra} & True \\
63 & \texttt{\texorpdfstring{\ensuremath{\ell}}{ℓ} = e[4] doubles ℍ: e[4] e[k] == e[4+k] (k = 1,2,3) and e[4] e[4] == -e[0]} & True \\
64 & \texttt{orientation inherited from 4+4.nb: e[1] e[2] == -e[3]} & True \\
65 & \texttt{Cayley-Dickson, γ = -1: (ac - d b̄, ā d + c b) reproduces the octonion product (symbolic)} & True \\
66 & \texttt{Cayley-Dickson, γ = +1: (ac + d b̄, ā d + c b) reproduces the SPLIT product of 4+4.nb (symbolic)} & True \\
67 & \texttt{the formula (ac - d̄ b, ā d + c b) of [1A] is NOT alternative: x·x·y != x(xy) for x = e[0]+e[4], y = e[1]} & True \\
68 & \texttt{... and its norm is not multiplicative (symbolic)} & True \\*
\end{longtable}}
\Needspace{6\baselineskip}\subsubsection*{3.6 Octonions versus split octonions}
{\footnotesize
\begin{longtable}{>{\raggedright\arraybackslash}p{\dimexpr 0.0700\linewidth-2\tabcolsep\relax}>{\raggedright\arraybackslash}p{\dimexpr 0.8100\linewidth-2\tabcolsep\relax}>{\raggedright\arraybackslash}p{\dimexpr 0.1200\linewidth-2\tabcolsep\relax}}
\toprule
\textbf{\#} & \textbf{check} & \textbf{result} \\
\midrule\endfirsthead
\toprule
\textbf{\#} & \textbf{check} & \textbf{result} \\
\midrule\endhead
\endfoot
\bottomrule\endlastfoot
69 & \texttt{e[A]e[B] (octonion) == s\_AB e[A]e[B] (split) for all A, B = 0..7} & True \\
70 & \texttt{over ℂ: ϕ(x *split y) == ϕ(x) *oct ϕ(y) (symbolic): split ⊗ ℂ ≅ \texorpdfstring{\ensuremath{\mathbb{O}}}{𝕆} ⊗ ℂ} & True \\
71 & \texttt{the norms correspond: N\_oct(ϕ(x)) == N\_split(x) (symbolic)} & True \\*
\end{longtable}}
\Needspace{6\baselineskip}\subsubsection*{4.1 The exact mirror of T16: the Euclidean Clifford algebra Cl(8,0)}
{\footnotesize
\begin{longtable}{>{\raggedright\arraybackslash}p{\dimexpr 0.0700\linewidth-2\tabcolsep\relax}>{\raggedright\arraybackslash}p{\dimexpr 0.8100\linewidth-2\tabcolsep\relax}>{\raggedright\arraybackslash}p{\dimexpr 0.1200\linewidth-2\tabcolsep\relax}}
\toprule
\textbf{\#} & \textbf{check} & \textbf{result} \\
\midrule\endfirsthead
\toprule
\textbf{\#} & \textbf{check} & \textbf{result} \\
\midrule\endhead
\endfoot
\bottomrule\endlastfoot
72 & \texttt{Cl(8,0): (T16O[A].T16O[B] + T16O[B].T16O[A])/2 == gO[[A+1,B+1]] ID16} & True \\
73 & \texttt{S44 == diag(σ, σ)} & True \\
74 & \texttt{T16O from the 4+4 matrices: T16O[A] == T16[A] (A = 0..3), T16O[k] == T16[k].S44 (k = 4..7)} & True \\
75 & \texttt{T16O[8] == T16[8] == diag(-ID8, ID8): same chirality matrix and same PL44, PR44 as 4+4.nb} & True \\*
\end{longtable}}
\Needspace{6\baselineskip}\subsubsection*{4.2 The REAL Clifford algebra of 7+1 spacetime: \texorpdfstring{\ensuremath{\Gamma}}{Γ}17, Cl(1,7) \texorpdfstring{\ensuremath{\cong}}{≅} \texorpdfstring{\ensuremath{\mathbb{R}}}{ℝ}(16)}
{\footnotesize
\begin{longtable}{>{\raggedright\arraybackslash}p{\dimexpr 0.0700\linewidth-2\tabcolsep\relax}>{\raggedright\arraybackslash}p{\dimexpr 0.8100\linewidth-2\tabcolsep\relax}>{\raggedright\arraybackslash}p{\dimexpr 0.1200\linewidth-2\tabcolsep\relax}}
\toprule
\textbf{\#} & \textbf{check} & \textbf{result} \\
\midrule\endfirsthead
\toprule
\textbf{\#} & \textbf{check} & \textbf{result} \\
\midrule\endhead
\endfoot
\bottomrule\endlastfoot
76 & \texttt{7+1 REAL Clifford relation: (Γ17[A].Γ17[B] + Γ17[B].Γ17[A])/2 == η17[[A+1,B+1]] ID16 == -η71[[A+1,B+1]] ID16} & True \\
77 & \texttt{Γ17[A] are REAL: every entry is 0 or ±1} & True \\
78 & \texttt{block form: Γ17[0] == \{\{0,ID8\},\{ID8,0\}\}, Γ17[i] == \{\{τO[i],0\},\{0,τObar[i]\}\}} & True \\
79 & \texttt{built from the real 4+4 matrices: Γ17[0] == T16[0], Γ17[i] == T16[0].T16[i] (i = 1,2,3), Γ17[k] == T16[0].T16[k].S44 (k = 4..7)} & True \\
80 & \texttt{Γ17[0] symmetric, Γ17[i] antisymmetric, so Transpose[Γ17[A]] == Γ17[0].Γ17[A].Γ17[0] (real Dirac adjoint)} & True \\
81 & \texttt{the 256 products Γ17[A1]...Γ17[Ak] (A1 < ... < Ak) are linearly independent: Cl(1,7) ≅ ℝ(16)} & True \\
82 & \texttt{grades: binomial(8,k) == \{1,8,28,56,70,56,28,8,1\}, total 256 (binomial(8,4) == 70, not 72)} & True \\
83 & \texttt{Tr[product] == 0 for every non-empty product, so Tr[M]/16 is the scalar part of M} & True \\
84 & \texttt{irreducible over ℝ: only multiples of ID16 commute with all Γ17[A] (Schur; real commutant of dimension 1)} & True \\
85 & \texttt{universal property: Γ17(v).Γ17(v) == (v.η17.v) ID16 for a symbolic vector v = X} & True \\*
\end{longtable}}
\Needspace{6\baselineskip}\subsubsection*{4.3 The complexified 7+1 realization in the exact block form of 4+4.nb: \texorpdfstring{\ensuremath{\Gamma}}{Γ}71}
{\footnotesize
\begin{longtable}{>{\raggedright\arraybackslash}p{\dimexpr 0.0700\linewidth-2\tabcolsep\relax}>{\raggedright\arraybackslash}p{\dimexpr 0.8100\linewidth-2\tabcolsep\relax}>{\raggedright\arraybackslash}p{\dimexpr 0.1200\linewidth-2\tabcolsep\relax}}
\toprule
\textbf{\#} & \textbf{check} & \textbf{result} \\
\midrule\endfirsthead
\toprule
\textbf{\#} & \textbf{check} & \textbf{result} \\
\midrule\endhead
\endfoot
\bottomrule\endlastfoot
86 & \texttt{7+1 COMPLEX Clifford relation, 4+4.nb sign: (Γ71[A].Γ71[B] + Γ71[B].Γ71[A])/2 == η71[[A+1,B+1]] ID16} & True \\
87 & \texttt{Γ71[0] == I T16[0] and Γ71[k] == T16O[k] (k = 1..7): only the time component is complex} & True \\
88 & \texttt{S = (ID16 - I Γ17[0])/Sqrt[2] is unitary and Γ71[A] == S.(I Γ17[A]).Inverse[S]} & True \\*
\end{longtable}}
\Needspace{6\baselineskip}\subsubsection*{4.4 Theorem: with \texorpdfstring{\ensuremath{\eta}}{η}71 there is no real 16×16 realization, because Cl(7,1) \texorpdfstring{\ensuremath{\cong}}{≅} \texorpdfstring{\ensuremath{\mathbb{H}}}{ℍ}(8)}
{\footnotesize
\begin{longtable}{>{\raggedright\arraybackslash}p{\dimexpr 0.0700\linewidth-2\tabcolsep\relax}>{\raggedright\arraybackslash}p{\dimexpr 0.8100\linewidth-2\tabcolsep\relax}>{\raggedright\arraybackslash}p{\dimexpr 0.1200\linewidth-2\tabcolsep\relax}}
\toprule
\textbf{\#} & \textbf{check} & \textbf{result} \\
\midrule\endfirsthead
\toprule
\textbf{\#} & \textbf{check} & \textbf{result} \\
\midrule\endhead
\endfoot
\bottomrule\endlastfoot
89 & \texttt{Schur: the intertwiners M with M.Conjugate[Γ71[A]] == Γ71[A].M form a 1-dimensional space} & True \\
90 & \texttt{... spanned by the volume element Γ17[8] == \{\{0,-ID8\},\{ID8,0\}\}, and Γ17[8].Γ17[8] == -ID16} & True \\
91 & \texttt{indicator(Γ71) == -1: Cl(7,1) ≅ ℍ(8), NO real 16×16 realization of the η71 relation} & True \\
92 & \texttt{indicator(Γ17) == +1: Cl(1,7) ≅ ℝ(16), real realization exists (Γ17)} & True \\*
\end{longtable}}
\Needspace{6\baselineskip}\subsubsection*{4.5 A real realization of Cl(7,1) needs 32×32 matrices}
{\footnotesize
\begin{longtable}{>{\raggedright\arraybackslash}p{\dimexpr 0.0700\linewidth-2\tabcolsep\relax}>{\raggedright\arraybackslash}p{\dimexpr 0.8100\linewidth-2\tabcolsep\relax}>{\raggedright\arraybackslash}p{\dimexpr 0.1200\linewidth-2\tabcolsep\relax}}
\toprule
\textbf{\#} & \textbf{check} & \textbf{result} \\
\midrule\endfirsthead
\toprule
\textbf{\#} & \textbf{check} & \textbf{result} \\
\midrule\endhead
\endfoot
\bottomrule\endlastfoot
93 & \texttt{REAL 32×32 realization of the η71 relation: (Γ71R32[A].Γ71R32[B] + Γ71R32[B].Γ71R32[A])/2 == η71[[A+1,B+1]] ID32} & True \\
94 & \texttt{Qi, Qj, Qk commute with every Γ71R32[A]; Qi.Qi == Qj.Qj == Qk.Qk == -ID32; Qi.Qj == Qk == -Qj.Qi} & True \\
95 & \texttt{the commutant of \{Γ71R32[A]\} has dimension exactly 4 and is spanned by ID32, Qi, Qj, Qk: it is ℍ} & True \\*
\end{longtable}}
\Needspace{6\baselineskip}\subsubsection*{4.6 Volume elements and chirality}
{\footnotesize
\begin{longtable}{>{\raggedright\arraybackslash}p{\dimexpr 0.0700\linewidth-2\tabcolsep\relax}>{\raggedright\arraybackslash}p{\dimexpr 0.8100\linewidth-2\tabcolsep\relax}>{\raggedright\arraybackslash}p{\dimexpr 0.1200\linewidth-2\tabcolsep\relax}}
\toprule
\textbf{\#} & \textbf{check} & \textbf{result} \\
\midrule\endfirsthead
\toprule
\textbf{\#} & \textbf{check} & \textbf{result} \\
\midrule\endhead
\endfoot
\bottomrule\endlastfoot
96 & \texttt{Γ17[8].Γ17[8] == -ID16 and Γ17[8] anticommutes with every Γ17[A]: no real chirality in 7+1} & True \\
97 & \texttt{Γ71[8].Γ71[8] == -ID16 and χ71 = -I Γ71[8] == T16O[8] == T16[8] (real, diagonal)} & True \\
98 & \texttt{7+1 chiral projectors in the Γ71 basis are exactly PL44, PR44 of 4+4.nb; χ71 anticommutes with every Γ71[A]} & True \\
99 & \texttt{complex chiral projectors for Γ17: PL17 + PR17 == ID16, PL17.PL17 == PL17, PL17.PR17 == 0, rank 8 each} & True \\*
\end{longtable}}
\Needspace{6\baselineskip}\subsubsection*{4.7 What the octonions do: \texorpdfstring{\ensuremath{\Gamma}}{Γ}(x)² = Re(x x), and alternativity}
{\footnotesize
\begin{longtable}{>{\raggedright\arraybackslash}p{\dimexpr 0.0700\linewidth-2\tabcolsep\relax}>{\raggedright\arraybackslash}p{\dimexpr 0.8100\linewidth-2\tabcolsep\relax}>{\raggedright\arraybackslash}p{\dimexpr 0.1200\linewidth-2\tabcolsep\relax}}
\toprule
\textbf{\#} & \textbf{check} & \textbf{result} \\
\midrule\endfirsthead
\toprule
\textbf{\#} & \textbf{check} & \textbf{result} \\
\midrule\endhead
\endfoot
\bottomrule\endlastfoot
100 & \texttt{Re(x x) == X.η17.X: the octonion square produces the 7+1 interval (symbolic)} & True \\
101 & \texttt{octonion basis: Γ17e[0] == \{\{0,ID8\},\{ID8,0\}\} and Γ17e[i] == \{\{LO[i],0\},\{0,-LO[i]\}\}} & True \\
102 & \texttt{alternativity in action: xImag(xImag ψ) == -|xImag|\textasciicircum{}2 ψ for the imaginary part xImag of x (symbolic)} & True \\
103 & \texttt{T16O(x)\textasciicircum{}2 == N(x) ID16 (octonion norm) and T16(x)\textasciicircum{}2 == (X.η4488.X) ID16 (4+4.nb)} & True \\*
\end{longtable}}
\Needspace{6\baselineskip}\subsubsection*{5.1 The Lorentz algebra so(1,7)}
{\footnotesize
\begin{longtable}{>{\raggedright\arraybackslash}p{\dimexpr 0.0700\linewidth-2\tabcolsep\relax}>{\raggedright\arraybackslash}p{\dimexpr 0.8100\linewidth-2\tabcolsep\relax}>{\raggedright\arraybackslash}p{\dimexpr 0.1200\linewidth-2\tabcolsep\relax}}
\toprule
\textbf{\#} & \textbf{check} & \textbf{result} \\
\midrule\endfirsthead
\toprule
\textbf{\#} & \textbf{check} & \textbf{result} \\
\midrule\endhead
\endfoot
\bottomrule\endlastfoot
104 & \texttt{28 real generators Σ17[A,B] (A < B), linearly independent} & True \\
105 & \texttt{[Σ17[A,B], Γ17[K]] == η17[[B+1,K+1]] Γ17[A] - η17[[A+1,K+1]] Γ17[B]} & True \\
106 & \texttt{so(1,7): [Σ\_AB, Σ\_KL] == η\_BK Σ\_AL - η\_AK Σ\_BL - η\_BL Σ\_AK + η\_AL Σ\_BK} & True \\
107 & \texttt{the complex structure J = Γ17[8] commutes with all Σ17[A,B] (Spin(1,7) acts complex-linearly on ℝ\textasciicircum{}16 = ℂ\textasciicircum{}8)} & True \\*
\end{longtable}}
\Needspace{6\baselineskip}\subsubsection*{5.2 A boost along x1 and a rotation in the x4–x5 plane}
{\footnotesize
\begin{longtable}{>{\raggedright\arraybackslash}p{\dimexpr 0.0700\linewidth-2\tabcolsep\relax}>{\raggedright\arraybackslash}p{\dimexpr 0.8100\linewidth-2\tabcolsep\relax}>{\raggedright\arraybackslash}p{\dimexpr 0.1200\linewidth-2\tabcolsep\relax}}
\toprule
\textbf{\#} & \textbf{check} & \textbf{result} \\
\midrule\endfirsthead
\toprule
\textbf{\#} & \textbf{check} & \textbf{result} \\
\midrule\endhead
\endfoot
\bottomrule\endlastfoot
108 & \texttt{S(β) and R(θ) are the exponentials: dS/dβ == Σ17[0,1].S, S(0) == ID16, dR/dθ == Σ17[4,5].R, R(0) == ID16, and S.Sinv == R.Rinv == ID16} & True \\
109 & \texttt{S(β).Γ17[K].S(β)\textasciicircum{}-1 == Σ\_B Λboost[[B+1,K+1]] Γ17[B]; Λboost is the boost along x1} & True \\
110 & \texttt{R(θ).Γ17[K].R(θ)\textasciicircum{}-1 == Σ\_B Λrot[[B+1,K+1]] Γ17[B]; Λrot is the rotation in the x4-x5 plane} & True \\
111 & \texttt{both Λ preserve the 7+1 metric: Transpose[Λ].η17.Λ == η17} & True \\*
\end{longtable}}
\Needspace{6\baselineskip}\subsubsection*{5.3 The 7+1 Dirac operator squares to the 7+1 wave operator}
{\footnotesize
\begin{longtable}{>{\raggedright\arraybackslash}p{\dimexpr 0.0700\linewidth-2\tabcolsep\relax}>{\raggedright\arraybackslash}p{\dimexpr 0.8100\linewidth-2\tabcolsep\relax}>{\raggedright\arraybackslash}p{\dimexpr 0.1200\linewidth-2\tabcolsep\relax}}
\toprule
\textbf{\#} & \textbf{check} & \textbf{result} \\
\midrule\endfirsthead
\toprule
\textbf{\#} & \textbf{check} & \textbf{result} \\
\midrule\endhead
\endfoot
\bottomrule\endlastfoot
112 & \texttt{(Γ17\textasciicircum{}A ∂\_A)\textasciicircum{}2 Ψ == (∂0\textasciicircum{}2 - ∂1\textasciicircum{}2 - ... - ∂7\textasciicircum{}2) Ψ for 16 arbitrary functions of x0..x7} & True \\*
\end{longtable}}
\Needspace{6\baselineskip}\subsubsection*{5.4 Spinors on the mass shell: real amplitudes}
{\footnotesize
\begin{longtable}{>{\raggedright\arraybackslash}p{\dimexpr 0.0700\linewidth-2\tabcolsep\relax}>{\raggedright\arraybackslash}p{\dimexpr 0.8100\linewidth-2\tabcolsep\relax}>{\raggedright\arraybackslash}p{\dimexpr 0.1200\linewidth-2\tabcolsep\relax}}
\toprule
\textbf{\#} & \textbf{check} & \textbf{result} \\
\midrule\endfirsthead
\toprule
\textbf{\#} & \textbf{check} & \textbf{result} \\
\midrule\endhead
\endfoot
\bottomrule\endlastfoot
113 & \texttt{p is on the mass shell (p.η17.p == m\textasciicircum{}2) and pslash.pslash == m\textasciicircum{}2 ID16} & True \\
114 & \texttt{(pslash - m) w == 0 has 8 linearly independent REAL solutions w} & True \\
115 & \texttt{Ψ = w Exp[-I p.x] solves (I Γ17\textasciicircum{}A ∂\_A - m) Ψ == 0} & True \\*
\end{longtable}}
\Needspace{6\baselineskip}\subsubsection*{5.5 Real spinor fields: massless with \texorpdfstring{\ensuremath{\Gamma}}{Γ}17, massive with the real 32×32 matrices}
{\footnotesize
\begin{longtable}{>{\raggedright\arraybackslash}p{\dimexpr 0.0700\linewidth-2\tabcolsep\relax}>{\raggedright\arraybackslash}p{\dimexpr 0.8100\linewidth-2\tabcolsep\relax}>{\raggedright\arraybackslash}p{\dimexpr 0.1200\linewidth-2\tabcolsep\relax}}
\toprule
\textbf{\#} & \textbf{check} & \textbf{result} \\
\midrule\endfirsthead
\toprule
\textbf{\#} & \textbf{check} & \textbf{result} \\
\midrule\endhead
\endfoot
\bottomrule\endlastfoot
116 & \texttt{the matrices anticommuting with every Γ17[A] are exactly the multiples of Γ17[8], and (c Γ17[8])\textasciicircum{}2 == -c\textasciicircum{}2 ID16: no real J with J.J == +ID16} & True \\
117 & \texttt{null momentum k = (1,1,0,0,0,0,0,0): k.η17.k == 0, kslash.kslash == 0, and kslash.w == 0 has 8 REAL solutions} & True \\
118 & \texttt{Ψ = w Cos[k.x] is a REAL solution of the real massless Dirac equation Γ17\textasciicircum{}A ∂\_A Ψ == 0} & True \\
119 & \texttt{(Γ71R32\textasciicircum{}A ∂\_A)\textasciicircum{}2 Ψ == -(∂0\textasciicircum{}2 - ∂1\textasciicircum{}2 - ... - ∂7\textasciicircum{}2) Ψ for 32 arbitrary functions of x0..x7} & True \\
120 & \texttt{p.η71.p == -m\textasciicircum{}2 and pslash32.pslash32 == -m\textasciicircum{}2 ID32} & True \\
121 & \texttt{Ψ = w Cos[p.x] + (pslash32.w/m) Sin[p.x], with 32 arbitrary REAL constants w, solves the REAL massive equation (Γ71R32\textasciicircum{}A ∂\_A - m) Ψ == 0} & True \\
122 & \texttt{... and Ψ obeys the Klein-Gordon equation (∂0\textasciicircum{}2 - ∂1\textasciicircum{}2 - ... - ∂7\textasciicircum{}2 + m\textasciicircum{}2) Ψ == 0} & True \\*
\end{longtable}}
\Needspace{6\baselineskip}\subsubsection*{5.6 All nine signatures p + q = 8, and the corrected table of [1B]}
{\footnotesize
\begin{longtable}{>{\raggedright\arraybackslash}p{\dimexpr 0.0700\linewidth-2\tabcolsep\relax}>{\raggedright\arraybackslash}p{\dimexpr 0.8100\linewidth-2\tabcolsep\relax}>{\raggedright\arraybackslash}p{\dimexpr 0.1200\linewidth-2\tabcolsep\relax}}
\toprule
\textbf{\#} & \textbf{check} & \textbf{result} \\
\midrule\endfirsthead
\toprule
\textbf{\#} & \textbf{check} & \textbf{result} \\
\midrule\endhead
\endfoot
\bottomrule\endlastfoot
123 & \texttt{split analogue: Re(x x) of split octonions == X.η53.X, and Γ53 realizes Cl(5,3) with REAL 16×16 matrices} & True \\
124 & \texttt{Γ08[A] = T16O[8].T16O[A] realizes Cl(0,8) with REAL 16×16 matrices} & True \\
125 & \texttt{gens[p] satisfy the Clifford relation of signature (p, 8-p) for p = 0..8} & True \\
126 & \texttt{classification: ℝ(16) for (8,0),(5,3),(4,4),(1,7),(0,8); ℍ(8) for (7,1),(6,2),(3,5),(2,6)} & True \\
127 & \texttt{every ℝ(16) case has explicit REAL 16×16 generators in this notebook (T16O, Γ53, T16, Γ17, Γ08)} & True \\
128 & \texttt{ℂ(16) cannot occur: its real dimension 2*16\textasciicircum{}2 = 512 differs from 2\textasciicircum{}8 = 256} & True \\*
\end{longtable}}
\Needspace{8\baselineskip}\subsection*{4.3 Independent Python cross-checks}\addcontentsline{toc}{subsection}{4.3 Independent Python cross-checks}
\textbf{tools/verify\_7p1\_nb.py (parses the evaluated 7+1.nb; exit code 0)}

\begin{Verbatim}[fontsize=\small,frame=single,framesep=2mm]
input cells                                                             : 56
output cells                                                            : 154
message cells (must be 0)                                               : 0
check outputs                                                           : 128
check outputs that are True                                             : 128
check["..."] calls in the input cells                                   : 128
$Version output                                                         : ['{"15.0.1 for
    ...Microsoft Windows (64-bit) (July 2, 2026)","Notebook$$43$40478`"}']
private notebook context                                                : True
Length[checks]                                                          : 128
failedChecks == {}                                                      : True
allChecksPassed output True                                             : True
tau[A] printed == 4+4 tau[A] (python)                                   : True
tauO[A] printed == python tauO[A]                                       : True
Gamma17[A] printed == python G17[A]                                     : True
gridO printed == python octonion table                                  : True
gridx44 printed == python split table (4+4.nb)                          : True
Lambda boost/rotation printed                                           : True
8 real on-shell spinors printed solve (pslash - m) w = 0 (python check) : True
classification table printed                                            : True

VERDICT: ALL CROSS-CHECKS PASSED
\end{Verbatim}
\textbf{tools/ref44.py (rebuilds 4+4.nb; exit code 0)}

\begin{Verbatim}[fontsize=\small,frame=single,framesep=2mm]
ζ7 = [-1, -1, -1, -1, 1, 1, 1, 1]
σ == ζ1ζ2ζ3 : True
{ζA,ζB}/2 == -η ID8 (A,B=1..7): True
-η_AA ζA == ζA^T : True
(τA τ̄B+τB τ̄A)/2 == η ID8 : True
(τ̄A τB+τ̄B τA)/2 == η ID8 : True
τ1..τ7 product == ID8: True
σ == τ4τ5τ6τ7: True
{T16A,T16B}/2 == η ID16 : True
T16[8] diag: [-1, -1, -1, -1, -1, -1, -1, -1, 1, 1, 1, 1, 1, 1, 1, 1]
σ16 == diag(-σ,σ): True
unit: [sqrt(2)/2, 0, 0, 0, sqrt(2)/2, 0, 0, 0]  unit.σ.unit = 1
E^a_A[h] == (u σ u) Transpose[Table[η_BB τ[B-1].u]] for all h: True
eA[1] : ['eA[1]', 'eA[2]', 'eA[3]', 'eA[4]', 'eA[5]', 'eA[6]', 'eA[7]', 'eA[8]']
eA[2] : ['eA[2]', '-eA[1]', '-eA[4]', 'eA[3]', '-eA[6]', 'eA[5]', 'eA[8]', '-eA[7]']
eA[3] : ['eA[3]', 'eA[4]', '-eA[1]', '-eA[2]', '-eA[7]', '-eA[8]', 'eA[5]', 'eA[6]']
eA[4] : ['eA[4]', '-eA[3]', 'eA[2]', '-eA[1]', '-eA[8]', 'eA[7]', '-eA[6]', 'eA[5]']
eA[5] : ['eA[5]', 'eA[6]', 'eA[7]', 'eA[8]', 'eA[1]', 'eA[2]', 'eA[3]', 'eA[4]']
eA[6] : ['eA[6]', '-eA[5]', 'eA[8]', '-eA[7]', '-eA[2]', 'eA[1]', '-eA[4]', 'eA[3]']
eA[7] : ['eA[7]', '-eA[8]', '-eA[5]', 'eA[6]', '-eA[3]', 'eA[4]', 'eA[1]', '-eA[2]']
eA[8] : ['eA[8]', 'eA[7]', '-eA[6]', '-eA[5]', '-eA[4]', '-eA[3]', 'eA[2]', 'eA[1]']
\end{Verbatim}
\textbf{tools/ref71.py (rebuilds 7+1.nb objects; exit code 0)}

\begin{Verbatim}[fontsize=\small,frame=single,framesep=2mm]
tauO antisymmetric (A>=1)                         : True
tauO[A]^2=-1 (A>=1)                               : True
tauO Euclid relation                              : True
tauObar tauO relation                             : True
tauO1..7 product                                  : [1, 1, 1, 1, 1, 1, 1, 1]
sigma commutes tau1..3, anticommutes tau4..7      : True
E orthogonal                                      : True
oct: unit two-sided                               : True
oct: e_i^2                                        : [-1, -1, -1, -1, -1, -1, -1]
oct: x xbar = N(x) 1                              : True
oct: N(xy)=N(x)N(y)                               : True
oct: conj(xy)=conj(y)conj(x)                      : True
oct: alternative (left,right,flexible)            : True
oct: Moufang identities                           : True
oct: #nonassociative basis triples                : 168
oct: #noncommuting ordered basis pairs            : 42
e_A*e_B (O) = s_AB e_A*e_B (S)                    : True
T16O Euclid Cl(8,0)                               : True
T16O[A]==T16A44[A] (A<=3)                         : True
T16O[k]==T16A44[k].S44 (k>=4)                     : True
S44==diag(sigma,sigma)                            : True
T16O[8] diag                                      : [-1, -1, -1, -1, -1, -1, -1, -1, 1, 1, 1,
    ...1, 1, 1, 1, 1]
G17 Cl(1,7)                                       : True
G17 real                                          : True
G17[0] sym, G17[i] antisym                        : True
G17[8]^2                                          : -1
G17 blocks: G17[i]==diag(tauO,tauObar)            : True
G71 Cl(7,1) eta71                                 : True
G71[8]^2                                          : -1
G71 = Ssim (I G17) Ssim^-1                        : True
G71R32 Cl(7,1) real                               : True
\end{Verbatim}
\textbf{tools/extra71.py (further identities; exit code 0)}

\begin{Verbatim}[fontsize=\small,frame=single,framesep=2mm]
Lambda boost (rows B = 0..7):
    ['cosh(beta)', '-sinh(beta)', '0', '0', '0', '0', '0', '0']
    ['-sinh(beta)', 'cosh(beta)', '0', '0', '0', '0', '0', '0']
    ['0', '0', '1', '0', '0', '0', '0', '0']
    ['0', '0', '0', '1', '0', '0', '0', '0']
    ['0', '0', '0', '0', '1', '0', '0', '0']
    ['0', '0', '0', '0', '0', '1', '0', '0']
    ['0', '0', '0', '0', '0', '0', '1', '0']
    ['0', '0', '0', '0', '0', '0', '0', '1']
Lambda rotation (rows B = 0..7):
    ['1', '0', '0', '0', '0', '0', '0', '0']
    ['0', '1', '0', '0', '0', '0', '0', '0']
    ['0', '0', '1', '0', '0', '0', '0', '0']
    ['0', '0', '0', '1', '0', '0', '0', '0']
    ['0', '0', '0', '0', 'cos(theta)', '-sin(theta)', '0', '0']
    ['0', '0', '0', '0', 'sin(theta)', 'cos(theta)', '0', '0']
    ['0', '0', '0', '0', '0', '0', '1', '0']
    ['0', '0', '0', '0', '0', '0', '0', '1']
indicator classification (signature, number of intertwiners, M.conj(M) for the normalized M,
    ...algebra):
    Cl(8,0): 1 intertwiner, M.conj(M) = +0.0625 * 1  ->  R(16)
    Cl(7,1): 1 intertwiner, M.conj(M) = -0.0625 * 1  ->  H(8)
    Cl(6,2): 1 intertwiner, M.conj(M) = -0.0625 * 1  ->  H(8)
    Cl(5,3): 1 intertwiner, M.conj(M) = +0.0625 * 1  ->  R(16)
    Cl(4,4): 1 intertwiner, M.conj(M) = +0.0625 * 1  ->  R(16)
    Cl(3,5): 1 intertwiner, M.conj(M) = -0.0625 * 1  ->  H(8)
    Cl(2,6): 1 intertwiner, M.conj(M) = -0.0625 * 1  ->  H(8)
    Cl(1,7): 1 intertwiner, M.conj(M) = +0.0625 * 1  ->  R(16)
    Cl(0,8): 1 intertwiner, M.conj(M) = +0.0625 * 1  ->  R(16)
EO == eta4488 . E44[8]                                                : True
G53 Cl(5,3) real                                                      : True
G08 Cl(0,8) real                                                      : True
so(1,7): [S_AB,S_CD]=eta_BC S_AD - eta_AC S_BD - eta_BD S_AC + eta_AD S_BC: True
[S_AB,G_C]=eta_BC G_A - eta_AC G_B                                    : True
S Sinv = 1                                                            : True
S G_C S^-1 == sum_B L[B,C] G_B                                        : True
L^T eta L == eta                                                      : True
p on shell p.eta17.p == 1                                             : True
pslash^2 == 1                                                         : True
8 real solutions                                                      : True
classification                                                        : True
explicit quaternion commutant of G71R32                               : True
G17e[0]==[[0,1],[1,0]]x1, G17e[i]==diag(LO[i],-LO[i])                 : True
Gamma17(x)^2 == (x.eta17.x) 1                                         : True
anticommutant of G17: dimension 1 (numerical rank 255 of 2048x256 system): True
G17[8] anticommutes with all G17[A], G17[8]^2 == -1                   : True
null k: k.eta17.k == 0, kslash^2 == 0, 8 rational solutions of kslash w = 0: True
Psi = w cos(k.x) is a REAL solution of G17^A d_A Psi = 0              : True
p.eta71.p == -m^2 and P32^2 == -m^2                                   : True
Psi32 (32 arbitrary real w) solves (G71R32^A d_A - m) Psi = 0         : True
Psi32 obeys the Klein-Gordon equation (box17 + m^2) Psi = 0           : True
(G71R32^A d_A)^2 == -box17 on 32 arbitrary functions                  : True
\end{Verbatim}
\Needspace{8\baselineskip}\subsection*{4.4 Hashes of the files used to build 7+1.nb and its documentation}\addcontentsline{toc}{subsection}{4.4 Hashes of the files used to build 7+1.nb and its documentation}
{\footnotesize
\begin{longtable}{>{\raggedright\arraybackslash}p{\dimexpr 0.3191\linewidth-2\tabcolsep\relax}>{\raggedright\arraybackslash}p{\dimexpr 0.6809\linewidth-2\tabcolsep\relax}}
\toprule
\textbf{file} & \textbf{SHA-256} \\
\midrule\endfirsthead
\toprule
\textbf{file} & \textbf{SHA-256} \\
\midrule\endhead
\endfoot
\bottomrule\endlastfoot
\texttt{7+1.nb} & \texttt{f7abe7279b5540e914de876dc055c199bd491b9c2d9522bf28dbc23689e237ce} \\
\texttt{4+4.nb} & \texttt{740023233eaa0cd24ca0e2e7010094671b1d314b698114c08db94050d3e0aed5} \\
\texttt{tools/cayley\_dickson\_formulas.py} & \texttt{bdde8c4c8199b09ef36ae4e15d5d53bc62744fc85b4c5ba52eab5516f77b5d5f} \\
\texttt{tools/cayley\_dickson\_match.py} & \texttt{73f2cc584bbe1a3620c50890ffe5fb299857b04a14bc34fed2f9eaefd2e36c4e} \\
\texttt{tools/classify\_clifford.py} & \texttt{42291276f46b046183fdbb7b62be7c5cedfd695fe9da96ac256afa23580930a9} \\
\texttt{tools/doc\_common.py} & \texttt{8db6983bb3afce5ecdceba0b1ce95ed3ce5427d3b3acbd67651cfbd931e46e5b} \\
\texttt{tools/doc\_provenance.py} & \texttt{12d4803e973f54837b9aa9aa92fb2bb0cc9a2342dc5cfad5439435174686d215} \\
\texttt{tools/doc\_readme.py} & \texttt{211c31b72adab6d1da0cf8cd5880b221cf36446ab099f0eea2e5be0cd96aa64d} \\
\texttt{tools/doc\_symbols.py} & \texttt{f7d7b363471d965d0e9fe900bebe6d10be700c3f02cb0684525fc9bbc5205b65} \\
\texttt{tools/docgen.py} & \texttt{bd229bca182d2de5939c5af0e9879e6ee22721629ee45518490b118661d76cfd} \\
\texttt{tools/extra71.py} & \texttt{9935ef99cd9300f2a7d4f06d36cf7e173912ebae61b6feff448ebe0588853168} \\
\texttt{tools/make\_7p1\_nb.py} & \texttt{6f617eb9d68adb69a8654529e5fb1d5287ee82e59ff48de96d373fcf6846c0ee} \\
\texttt{tools/make\_docs.py} & \texttt{b89ef1f433cf259099cc359a0bdf9b2edb7d0fc148211890a953105a6877c05f} \\
\texttt{tools/nbparse.py} & \texttt{ad2255a8864be7447a55023840b327f789cb8c9d868d42261b19261653847423} \\
\texttt{tools/nbwrite.py} & \texttt{dcbe8ce5b3203c89a1d17fbe946e1b6b7bb7f8895afa944743d651cfefadeffb} \\
\texttt{tools/octalg.py} & \texttt{36b50120b69a963e252e07d60ec007755220822b4ad339716ea7aed93d58eb24} \\
\texttt{tools/ref44.py} & \texttt{ab100c3ebcc055fbe59c132a02f80ecb17630e934e99d59f88e919ec8064faf5} \\
\texttt{tools/ref71.py} & \texttt{4ac783b1ab221317387f2a74e0be60f114fd64940fe3ea96af79b33ee7dcb8b0} \\
\texttt{tools/run\_eval.wl} & \texttt{f0f89ead42f002f839709b335be7007e0112c4d98f87a5768fd148bcac4de876} \\
\texttt{tools/verify\_7p1\_nb.py} & \texttt{0521c82e157d1712c48b51a572d3f8d28de58118aa7d639b4378fd39f31eb558} \\*
\end{longtable}}
\Needspace{12\baselineskip}\section*{5. Process and environment}\addcontentsline{toc}{section}{5. Process and environment}
\begin{itemize}
\item Workspace: \texttt{D:\textbackslash{}Developer\textbackslash{}github\textbackslash{}7+1} on Windows; Python and LaTeX ran in a Linux shell with that folder mounted. Intermediate files go to \texttt{build/}, which the repository's \texttt{.gitignore} excludes.
\item The notebook is generated by \texttt{tools/make\_7p1\_nb.py} (cells written with \texttt{tools/nbwrite.py}), evaluated in place by \texttt{tools/run\_eval.wl} in Wolfram 15.0.1, and cross-checked by \texttt{tools/verify\_7p1\_nb.py}.
\item Python 3.10.12 with sympy 1.14.0 and numpy 2.2.6 (exact rational arithmetic for every identity).
\item The documentation is generated by \texttt{tools/make\_docs.py} from \texttt{tools/docgen.py}, \texttt{tools/doc\_common.py}, \texttt{tools/doc\_readme.py}, \texttt{tools/doc\_symbols.py} and \texttt{tools/doc\_provenance.py}: one source renders Markdown and LaTeX, and every matrix, table and check list is computed or parsed, never typed. The LaTeX is compiled with XeLaTeX (TeX Live 2022, \texttt{latexmk -xelatex}, fonts Latin Modern and DejaVu Sans Mono).
\item History: the first algebraic experiments were made in a temporary cloud sandbox and then moved into \texttt{tools/} of the workspace, where everything was re-run. A Wolfram 15.0 kernel window was opened once by mistake before Wolfram 15.0.1 was available; no result comes from it. Only Wolfram 15.0.1 evaluated 7+1.nb.
\item Corrections made during the review: an earlier draft called the real on-shell amplitudes of notebook section 5.4 "Majorana spinors"; that label was removed, and section 5.5 was added to show which real spinor FIELDS exist (massless with 16 components, massive only with the 32×32 matrices).
\item The evaluated notebook was rendered section by section in the Wolfram 15.0.1 front end (\texttt{Rasterize}) and inspected. Characters that the front end's text font shows badly are written differently in text cells: the double-struck O and the arrow of $\vec x$ as small inline formulas, the script $\ell$ in the font Cambria Math, and $\not p$, $\not k$ as \texttt{pslash}, \texttt{kslash}.
\end{itemize}
\Needspace{12\baselineskip}\section*{6. Authorship}\addcontentsline{toc}{section}{6. Authorship}
Notebook, tools and documentation written by Claude (Anthropic, model claude-opus-5-5) for Nash, on 2026-10-26, from the task statement and the notebook 4+4.nb described in section 1.

\end{document}
